% Part: lambda-calculus
% Chapter: lambda-definability
% Section: pairs

\documentclass[../../../include/open-logic-section]{subfiles}

\begin{document}

\olfileid{lam}{ldf}{pai}
\olsection{క్రమయుగ్మాలు మరియు పూర్వవర్తి}

\begin{defn}
$M$, $N$ అనే పదాల క్రమయుగ్మాన్ని ($\tuple{M,N}$గా రాస్తాం) ఇలా
నిర్వచిస్తాం:
\[
\tuple{M,N} \ident \lambd[f][fMN].
\]
\end{defn}

సహజంగా చూస్తే, ఇది ఒక ప్రమేయాన్ని స్వీకరించి, దాన్ని
క్రమయుగ్మంలోని రెండు అవయవాలకు ప్రయోగించే ప్రమేయం. ఈ ఆలోచనను
అనుసరించి రెండు పదాలను స్వీకరించి, వాటితో కూడిన క్రమయుగ్మాన్ని
తిరిగి ఇచ్చే నిర్మాణ ప్రమేయం ఇది:
\[
  \fn{Pair} \ident \lambd[mn][\lambd[f][fmn]]
\]
క్రమయుగ్మం ఇచ్చినప్పుడు దాని అవయవాలను కూడా తిరిగి పొందాలనుకుంటాం.
అందుకు క్రమయుగ్మాన్ని ఆర్గ్యుమెంటుగా స్వీకరించి, దాని మొదటి లేదా
రెండవ అవయవాన్ని తిరిగి ఇచ్చే రెండు ఎంపిక ప్రమేయాలు కావాలి:
\begin{align*}
  \fn{Fst} & \ident \lambd[p][p(\lambd[mn][m])]\\
  \fn{Snd} & \ident \lambd[p][p(\lambd[mn][n])]
\end{align*}

\begin{prob}
  $\fn{Fst}$, $\fn{Snd}$ అనే ఎంపిక ప్రమేయాలు ఎందుకు
  పనిచేస్తాయో వివరించండి.
\end{prob}

ఇప్పుడు క్రమయుగ్మాల సాయంతో పూర్వవర్తి ప్రమేయాన్ని
\tetoken{లాంబ్డాతో నిర్వచించవచ్చు}{!!{lambda define}}:
\[
  \fn{Pred} \ident \lambd[n][\fn{Fst}(n (\lambd[p][\tuple{\fn{Snd}\, {p}, \fn{Succ}(\fn{Snd}\, {p})}]) \tuple{\num 0, \num 0})]
\]
$\num n\, f x$ అనేది $f^{n}(x)$కు తగ్గుతుందని గుర్తుంచుకోండి.
ఇక్కడ $f$ ఒక క్రమయుగ్మం~$p$ను స్వీకరించి, $p$ రెండవ అవయవం,
ఆ అవయవపు ఉత్తరగామితో కూడిన కొత్త క్రమయుగ్మాన్ని తిరిగి ఇచ్చే ప్రమేయం;
$x$ అనేది $\tuple{0,0}$ అనే క్రమయుగ్మం. కాబట్టి $n=0$ అయితే
ఫలితం $\tuple{0,0}$; లేకపోతే
$\tuple{\num{n-1}, \num n}$. ఆ ఫలితపు మొదటి
అవయవాన్ని $\fn{Pred}$ తిరిగి ఇస్తుంది.

తీసివేతను $\fn{Pred}$ను $a$కు $b$ సార్లు ప్రయోగించడం
ద్వారా నిర్వచించవచ్చు:
\[
\fn{Sub} \ident \lambd[ab][b \fn{Pred}\, a].
\]

\end{document}
